\documentclass[letterpaper,10pt]{article}

\IfFileExists{lmodern.sty}{\usepackage[T1]{fontenc}\usepackage{lmodern}}{}
\usepackage[left=0.55in,right=0.55in,top=0.45in,bottom=0.45in]{geometry}
\usepackage{hyperref}
\usepackage{titlesec}
\usepackage{enumitem}
\usepackage{tabularx}
\usepackage[expansion=false]{microtype}
\DeclareTextSymbolDefault{\textbullet}{OMS}

\hypersetup{colorlinks=true,urlcolor=blue,linkcolor=blue}
\urlstyle{same}
\pagestyle{empty}
\raggedbottom
\setlength{\parindent}{0pt}
\setlength{\parskip}{2.2pt}
\setlist[itemize]{leftmargin=1.3em,itemsep=1.3pt,topsep=1.5pt,parsep=0pt,partopsep=0pt}
\titleformat{\section}{\normalsize\bfseries\uppercase}{}{0em}{}[\vspace{-2pt}\titlerule]
\titlespacing*{\section}{0pt}{9pt}{4.5pt}

% \entry{left, bold}{right}
\newcommand{\entry}[2]{%
  \noindent\textbf{#1}\hfill #2\par\nopagebreak\vspace{-0.5pt}
}
\newcommand{\proj}[1]{\vspace{1.5pt}\textbf{\textit{#1}}}


\begin{document}

\begin{center}
    {\LARGE\textbf{Andrei Medvedev}}\\[3pt]
    Seattle, WA
    \quad $|$ \quad
    \href{mailto:am334@uw.edu}{am334@uw.edu}
    \quad $|$ \quad
    \href{https://github.com/mdvdvandrei}{github.com/mdvdvandrei}
\end{center}
\vspace{-6pt}

\textbf{Research interests:} ocean circulation and climate variability, generative ocean and climate modeling, data assimilation, probabilistic reconstruction from sparse observations

% =====================================================================
\section{Education}

\entry{University of Washington, Seattle, WA}{2024 -- June 2029 (expected)}
PhD, Physical Oceanography, generative machine learning for climate and ocean state estimation. GPA 3.95/4.00\\
Advisors: Georgy Manucharyan, Kyle Armour. NASA FINESST Fellowship (\$150,000, about 13\% selection rate)

\entry{Lomonosov Moscow State University, Moscow, Russia}{2018 -- 2024}
MS, Physics (2024), GPA 3.92/4.00, thesis on Bayesian inference of river runoff change from a multi-model climate ensemble\\
BS, Physics and Applied Mathematics (2022), thesis on deep reinforcement learning for autonomous pursuit and interception\\
State Scholarship (top 10\% of students)

% =====================================================================
\section{Research Experience}

\entry{Graduate Research Assistant, University of Washington}{2024 -- Present}

\proj{Generative data assimilation from sparse observations} (first author; \href{https://doi.org/10.22541/essoar.15011822/v1}{preprint}, \href{https://github.com/mdvdvandrei/GenDA-historical-3d-ocean-reconstruction}{code})
\begin{itemize}
    \item Trained a video diffusion model (EDM, cBottle architecture, 73.5M parameters) from scratch on a 50-member climate-model ensemble as a prior over coupled ocean-atmosphere trajectories on a HEALPix grid.
    \item Sampled posteriors zero-shot via likelihood guidance from sparse, irregular observations (EN4 ocean temperature and salinity profiles), denoising all 852 months of 1950 to 2020 jointly across GPUs.
    \item Evaluated on held-out members, a different climate model and the real ocean. Unassimilated sea level tracks satellite altimetry, and rank histograms trace ensemble underdispersion to the sampler.
\end{itemize}

\proj{AMOC changes since 1900 from sea-surface observations} (first author, \textit{Nature Geoscience}, under review)
\begin{itemize}
    \item Built heteroscedastic residual networks predicting Atlantic overturning strength and its variance from surface fields. Trained only on 18 climate models, they reach ${\mathbf{r = 0.82}}$ on unseen direct measurements, above reanalyses.
    \item Showed that training on 18 simulators vs.\ 1 drives transfer to real data, and used Shapley values to tie predictions to physically meaningful regions. Found no 20th-century decline and a weakening only after 2000.
\end{itemize}

\proj{Additional projects}
\begin{itemize}
    \item Extended the NeurOST sea-level product (deep learning on satellite altimetry and SST) to a 2003 to 2025 record.
    \item Trained models on \textbf{up to 32 GPUs} (NASA A100 and GH200 nodes) with PyTorch DDP, on \textbf{10+ TB} of simulations and satellite data.
\end{itemize}

\vspace{2pt}
\entry{Research Assistant, A.M. Obukhov Institute of Atmospheric Physics, Moscow}{2022 -- 2024}

% =====================================================================
\section{Industry Experience}

\entry{Machine Learning Intern (NLP), SkillCup \textnormal{(edtech startup)}}{2023}
\begin{itemize}
    \item Distilled a larger LLM into FRED-T5-large (820M-parameter transformer) for quiz generation from course texts and lecture transcripts, via synthetic pairs from the larger model plus a curated question bank, deployed with FastAPI.
\end{itemize}

% =====================================================================
\section{Publications and Presentations}

{\small
\textbf{A. Medvedev}, H. Wei, K. Armour, K. Srinivasan, A. Solodoch, A. Stewart, G. Manucharyan. Basin-wide sea-surface observations indicate recent AMOC weakening without a 20th-century decline. \textit{Nature Geoscience}, under review. Preprint on \href{https://doi.org/10.22541/essoar.15002298/v1}{ESS Open Archive}.

\textbf{A. Medvedev}, V. Cooper, G. Manucharyan. Reconstructing the global ocean from the historical observing record with generative data assimilation. Preprint on \href{https://doi.org/10.22541/essoar.15011822/v1}{ESS Open Archive}, to be submitted to \textit{JAMES}.

H. Wei, A. Stewart, \textbf{A. Medvedev}, K. Srinivasan, A. Solodoch, G. Manucharyan, A. Hogg. Deep learning unlocks satellite-based monitoring of the global ocean overturning circulation. Preprint on \href{https://essopenarchive.org/doi/abs/10.22541/essoar.15004679}{ESS Open Archive}.

\textbf{A. Medvedev}, A. Eliseev, I. Mokhov. Bayesian estimates for changes of the Russian river runoff in the 21st century as based on the CMIP6 model ensemble simulations. \textit{Izvestiya, Atmospheric and Oceanic Physics}, 2024.

Presentations: AGU Fall Meeting 2026 (two oral talks); AI4Ocean Workshop 2026 (lightning talk); AGU Fall Meeting 2025 (poster).\par}

% =====================================================================
\section{Skills}

\textbf{Generative models:} diffusion and video diffusion (EDM), flow matching, diffusion bridges, guided posterior sampling\\
\textbf{Machine learning:} transformers (fine-tuning, distillation), probabilistic scoring (NLL, CRPS, rank histograms)\\
\textbf{Environmental data:} in-situ ocean profiles, satellite altimetry and SST, CMIP6 ensembles, HEALPix grids\\
\textbf{Engineering:} Python, PyTorch (DDP), Hydra, Bash, Git, Docker, Apptainer, PBS/Slurm; Zarr, xarray, Dask

\end{document}
